% !TEX root = ../enac_project_fr.tex
\chapter*{Introduction générale}
\label{introduction}
\addcontentsline{toc}{chapter}{\nameref{introduction}}
\markright{Introduction générale}
\thispagestyle{fancy}

Les cursus de l'ENAC comprennent des stages et des projets de synthèse. Ces derniers doivent permettre aux étudiants de mettre en application les enseignements qu'ils ont reçus.\\
L'objectif de ce document est de fournir aux étudiants des recommandations à respecter impérativement pour la rédaction de leur rapport et pour la préparation de leur soutenance. Les conseils donnés ici sont généraux et pourront donc être repris lors de tout travail technique ou scientifique. Ce document constitue également un modèle pour ceux rédigeant en \LaTeX. \\
Le chapitre 1 concerne la rédaction du rapport tandis que le chapitre 2 aborde la soutenance.  

\section{Présentation du contexte}

todo 

\section{Problématique générale}

todo 

